\section{Two Morita transformations and an \texorpdfstring{$A_4$}{A4} obstruction}
\label{sec:c3-a4}
The following argument uses actual module categories of a hypothetical
categorification. Its finite checks only verify the source products,
eight alternating forms, and the elementary ring normalization.

\begin{theorem}
\label{thm:c3-a4}
The rings $R7\_2\_236$, $R7\_4\_2167$, and $R7\_4\_2716$
have no categorification over $\C$.
\end{theorem}
\begin{proof}
Their dimension vectors are $(1,1,1,3,d,d,d)$, with
$d^2=4+md$ and $m=6,12,12$, respectively. The invertibles form
$C_3=\langle g\rangle$, fixing the dimension-three simple $X$ and
permuting the other three simples freely. The source products give
$X^2=1+g+g^2+2X$.

\emph{The first Morita transformation.}
The fixed simple $X$ trivializes the pointed associator, so the
regular algebra $A=1+g+g^2$ exists. Since
$H^2(C_3,\C^\times)=0$, $X$ has both left and right $A$-module
structures. Thus $X\otimes X$ is an $A$-bimodule. Its $X$-isotypic
part is $2X$. A nontrivial projective class for $C_3\times C_3$
has irreducible degree three, so this multiplicity-two part forces
the class on $X$ to be trivial. It gives nine simple bimodules with
underlying object $X$, each of dual dimension one. The pointed
orbit gives three more invertibles. The remaining orbit has cyclic
stabilizer and gives three simples of dimension $d$.

Consequently $\mathcal C_A$ has dimensions $1^4,d$, and its Morita
dual $\mathcal D$ has dimensions $1^{12},d^3$. Already this excludes
$m=6$: the published \cite[Proposition~3.3]{Schopieray} requires
$s\mid r$ for a category with two simple dimensions $1,d$ and
irrational $d$ satisfying $d^2=s+rd$. Here $4\nmid6$.

For the two remaining cases, write $K$ for the twelve invertibles
of $\mathcal D$. They fix the unique $d$-dimensional object of
$\mathcal C_A$; its rank-one module structure trivializes the whole
pointed associator. Thus $\mathcal D_{\rm pt}\simeq\mathrm{Vec}_K$.
The integral subcategory of $\mathcal C$ has rank four, and its
$A$-module category has four simple objects on which $K$ acts
transitively. A stabilizer $L$ has order three. If $L$ were normal,
the dual of $\mathrm{Vec}_K$ on this module category would have
four double cosets, each contributing three characters of $L$,
hence rank twelve instead of four. Therefore $L$ is nonnormal. Its core is trivial because $|L|=3$,
so the action embeds $K$ as an index-two subgroup of $S_4$; hence
$K=A_4$.

For a noninvertible $Y$ of $\mathcal D$, comparison of rational and
irrational coefficients in $\FPdim(YY^*)=4+12d$ shows that $YY^*$
contains exactly four invertibles. Each occurs with multiplicity at
most one by rigidity. Thus both stabilizers are the unique subgroup
$H=V_4$ of $A_4$, and the three noninvertibles form one orbit.

\emph{The second Morita transformation.}
At least one of these three simples, say $Y$, is self-dual. The
projective $H\times H$ stabilizer of $Y$ has an alternating
commutator form on $\mathbb F_2^4$. Duality exchanges the left and
right factors and inverts the commutator. Its values are $\pm1$,
so in matched bases the form is
\[
 \beta=\begin{pmatrix}bJ&C\\ C&bJ\end{pmatrix},\qquad
 J=\begin{pmatrix}0&1\\1&0\end{pmatrix},\quad C=C^{\mathsf t}.
\]
This uses rigidity, with no pivotal assumption. The two algebra
structures on $\bigoplus_{h\in H}h$ differ by the unique nontrivial
class in $H^2(H,\C^\times)$, adding $\operatorname{diag}(J,J)$ to
$\beta$. For the eight symmetric matrices $C$, the two ranks are
$2,4$ in some order, except for $C=0$, when they are $0,4$.
The rank-zero alternative gives a simple bimodule of dimension
$d/4$. Its monic minimum polynomial has constant coefficient
$-1/4$, contradicting algebraic integrality. Hence one algebra
choice gives rank two. The corresponding twisted group algebra
has four irreducible representations of degree two, giving four
simple bimodules over each $Y$, of dimension $t=d/2$.
Transport by $A_4$ preserves the unique nontrivial $H$-cocycle
class, so the same degree holds over all three noninvertibles.

The pointed Morita dual is again $\mathrm{Vec}_{A_4}$ for either
algebra choice. For the untwisted algebra this is Fourier duality
for $V_4\rtimes C_3$. The other cocycle extends to $A_4$ through
$Q_8\rtimes C_3\to A_4$, so an identity-on-objects tensor
autoequivalence of the pointed subcategory identifies the two
pointed duals. No extension of this autoequivalence to $\mathcal D$
is assumed. The resulting category $\mathcal D'$ therefore has
dimensions $1^{12},t^{12}$, with invertible group $A_4$ and
$t^2=1+6t$. For every noninvertible $Z$, comparison of rational
coefficients in $\FPdim(ZZ^*)=1+6t$ gives a trivial stabilizer.
Thus these twelve objects form a free $A_4$-orbit.

\emph{Normalization without a self-duality assumption.}
Choose any noninvertible $\rho$ in $\mathcal D'$. Right
multiplication is free as well, since duality turns a right
stabilizer into a left stabilizer. Write
\[
 \rho g=\phi(g)\rho,\qquad \rho^2=a+N\rho,
 \qquad N\in\Z_{\geq0}[A_4],\quad \sum_g n_g=6.
\]
Rigidity gives a unique invertible summand $a$. Associativity gives
\[
 \phi^2=\operatorname{Ad}_a,\quad \phi(a)=a,\quad
 Ns=\phi(s)N,\quad \phi(N)=N,
\]
and then $\rho^*=a^{-1}\rho$ and $N^*=a^{-1}N$.
Since $\operatorname{Aut}(A_4)=S_4$ and $Z(A_4)=1$, write
$\phi=\operatorname{Ad}_\sigma$; necessarily $a=\sigma^2$.
Replacing $\rho$ by $b\rho$ replaces $(\sigma,a,N)$ by
$(b\sigma,(b\sigma)^2,bN)$. Choose $b=\sigma^{-1}$ when $\sigma$
is even, and $b=\tau\sigma^{-1}$ otherwise, with $\tau$ a
transposition. In either case $b\in A_4$, the new $a$ is one,
and the new $\rho$ is self-dual. The two normalized actions are
therefore the identity and conjugation by $\tau$.

For the identity action, centrality and self-adjointness give
\[
 N=A\,1+B\sum_{h\in V_4\setminus\{1\}}h,
 \qquad (A,B)=(6,0),(3,1),(0,2).
\]
The two size-four three-cycle classes would have equal
coefficients, so neither can occur in a sum of size six.
The first choice contains the forbidden rank-two subcategory
$\rho^2=1+6\rho$ \cite{OstrikRank2}.
For the other choices, $N$ acts on the three-dimensional
representation of $A_4$ by $A-B=\pm2$. Taking $\rho$ to act by
a root of $x^2\mp2x-1$ gives the formal codegrees
\[
                         16+8\sqrt2,\qquad16-8\sqrt2.
\]
For the transposition action, $\tau N$ is constant on the two
$A_4$-conjugacy classes in the odd coset of $S_4$: the six
transpositions and the six four-cycles. It must be one of their
class sums. On the standard three-dimensional representation of
$S_4$, these class sums act by $2$ and $-2$, respectively.
Taking $\rho=rT$, where $T$ represents $\tau$ and
$r^2=1\pm2r$, again gives the displayed formal codegrees.
In both computations the Casimir is $12(1+\rho^2)$ and its
scalar equals the representation degree times its formal codegree.

Finally, the published \cite[p.~1628, Lemma~3.1]{Schopieray}
states that every category with two simple Frobenius--Perron
dimensions is Galois conjugate to a pseudo-unitary category, with
the same based ring. In this conjugate, the central constituent
$A_E$ associated to an irreducible representation has dimension
$D/f_E$, where $D=\FPdim(\mathcal D')$
\cite[Theorem~2.13]{OstrikRank3}. Its forgetful image has dimension
in $\Z+\Z t$, so every $f_E$ must belong to $\mathbb Q(t)$.
Here $t=3+\sqrt{10}$, whereas $16\pm8\sqrt2$ do not belong to
$\mathbb Q(\sqrt{10})$. This contradiction excludes both remaining
source rings.
\end{proof}

For $R7\_3\_1022$, the same reduction has $m=8$ and $t^2=1+4t$.
Its normalized system retains the branch
$N=\sum_{h\in V_4}h$, which passes the formal-codegree field
condition. Thus this argument makes no categorifiability decision
for that source ring.
